% Generated draft: requires mathematical and source review.
\documentclass[11pt]{article}
\usepackage[margin=25mm]{geometry}
\usepackage{fontspec}
\setmainfont{Latin Modern Roman}
\usepackage{amsmath,amssymb,mathtools,unicode-math}
\setmathfont{Latin Modern Math}
\usepackage{enumitem,needspace,xurl,hyperref}
\setlength{\parindent}{0pt}
\setlength{\parskip}{5pt}
\setlength{\emergencystretch}{3em}
\newcommand{\N}{\mathbb{N}}
\newcommand{\Z}{\mathbb{Z}}
\newcommand{\Q}{\mathbb{Q}}
\newcommand{\R}{\mathbb{R}}
\newcommand{\C}{\mathbb{C}}
\newcommand{\F}{\mathbb{F}}
\newcommand{\A}{\mathbb{A}}
\newcommand{\PP}{\mathbb{P}}
\newcommand{\E}{\mathbb{E}}
\newcommand{\eps}{\varepsilon}
\newcommand{\dd}{\,\mathrm{d}}
\begin{document}
\section*{UnsolvedMath ID: 30002487}
\textbf{Source problem code:} catalogue entry
\subsection{Definitions and mathematical statement}
The setting is burnstile theory for propagating fronts. The supplied record refers to stable and unstable burning invariant manifolds, primary intersection points (abbreviated pips), and a sequence of modified entrance and exit lobes. It does not supply mathematical definitions of these objects, the governing front dynamics, or the class of systems under consideration. These terms therefore retain their source-specific meanings; the questions below require those definitions before they can become fully quantified mathematical statements.

\begin{enumerate}
\item Characterize the structure of the sequence of modified entrance and exit lobes. Is there a unique locally smallest lobe? What transport questions are answered by the other lobes?
\item Can a continuous-time treatment improve efficiency? The supplied record specifies neither the treatment to be compared nor a measure of efficiency.
\item What are the physical consequences when the stable and unstable burning invariant manifolds lose a primary intersection point, particularly for ballistic front propagation?
\item Can burnstile theory be extended continuously through the loss of a primary intersection point? The original wording asks whether the theory can accommodate the loss ``smoothly''; the record does not specify whether this means continuity, differentiability, or another regularity requirement.
\end{enumerate}

For the first question, the indexing and construction of the lobe sequence, the criterion for comparing lobe sizes, and the meaning of locality are unspecified. For the last two questions, the mechanism of intersection loss, the varying parameter or time variable, the relevant propagation observable, and the topology for an extension are unspecified. No particular choices of these missing data are asserted here.
\subsection{Short English statement}
The problem asks how modified entrance and exit lobes organize transport in burnstile theory, whether a continuous-time approach improves efficiency, and how losing a primary intersection affects front propagation and the continuation of the theory. The supplied record lacks the definitions and comparison criteria needed for a fully precise formulation.
\subsection{Sources}
\begin{itemize}
\item Local UnsolvedMath registry, permanent ID \(30002487\), exact source problem code \texttt{catalogue entry}, title: \emph{Burnstile Lobe Sequences and Loss of Primary Intersections}. Collection: Oberwolfach Reports Open Problems; proposed year recorded as 2014. No individual proposer is identified in the supplied record.
\item Supplied report citation: \emph{Mixing, Transport and Coherent Structures} (2014), Oberwolfach Reports, pp.\ 213--286, DOI \texttt{10.4171/owr/2014/04}. Supplied source URL and literature-source link, labelled ``OWR report'': \url{https://doi.org/10.4171/owr/2014/04}. The record describes these reports as documenting mathematical workshops and their open-problem discussions; it supplies no narrower page or passage for this problem.
\item Literature assessment dated \texttt{2026-08-21}: the supplied curation labels the problem open and reports that the cited report poses these research questions and that the literature located gives no general answer covering them. This is supplied dated curation, not an independently verified literature assessment. The background explicitly requests independent verification of time-sensitive status and cited claims.
\item The registry labels the statement \texttt{corrected\_verified} and reports that it was checked against the official report, with the clean statement restoring source hypotheses and notation while preserving the questions. This verification claim is supplied metadata; it has not been independently verified in this draft.
\end{itemize}
\end{document}
